\section*{Rangkuman}
\addcontentsline{toc}{section}{Rangkuman}

Pada bab ini telah dibahas:
\begin{itemize}[leftmargin=*]
  \item notasi himpunan, relasi, fungsi, dan induksi matematika;
  \item alfabet, string, bahasa, serta operasi konkatenasi dan
        bintang Kleene;
  \item tata bahasa sebagai sistem penulisan ulang, derivasi, dan
        pohon urai;
  \item hierarki Chomsky dan mesin pengenal yang bersesuaian;
  \item data pendukung tabel pemetaan Tipe 3--0 dan derivasi/pohon;
  \item peta konsep dari teori bahasa ke analisis leksikal,
        sintaksis, dan semantik.
\end{itemize}

Bab III akan memakai seluruh notasi ini untuk menjelaskan arsitektur
kompilator. Kita akan melihat bagaimana string sumber berubah menjadi
token, pohon, dan representasi semantik, serta di mana setiap teknik
yang dipelajari pada bab-bab berikutnya digunakan.
